%% ma: L11-B09-0008
%% lop: 11
%% chuong: 3
%% bai: 9
%% dang: 11.9.6
%% loai: TN4
%% muc: NB
%% dap_an: "D"
%% nguon: "de-on-tap-giua-ki-1-toan-11-nam-2026-2027-truong-thpt-ngo-quyen-da-nang.pdf (M09, Đề ôn tập giữa kì 1 Toán 11 THPT Ngô Quyền, Đà Nẵng), Mã đề 0111, Phần I Câu 8"
%% kien_thuc: "Bài 8 (mẫu số liệu ghép nhóm), Bài 9 (các số đặc trưng đo xu thế trung tâm)"
%% trang_thai: chinh-thuc
%% ghi_chu: "Đã duyệt 29/10/2026."
\begin{ex}
Cho mẫu số liệu ghép nhóm về điểm kiểm tra của học sinh. Giá trị đại diện của nhóm số liệu $[a;b)$ được tính bằng công thức
\choice{$b-a$}{$\dfrac{b-a}{2}$}{$a+b$}{\True $\dfrac{a+b}{2}$}
\loigiai{Giá trị đại diện của nhóm $[a;b)$ là trung điểm $\dfrac{a+b}2$.}
\end{ex}
